\section{Conserved internal correlations and the critical mass}
\label{sec:r12-internal-dynamics}

The quantitative polarization theorem has a consequence for every time
of an actual solution. The internal correlation operator is conserved,
and a field with more than one occupied spatial mode cannot approach
the scalar sharp interpolation constant arbitrarily closely. This
raises the mass range in which conservation of energy controls the
whole evolution.

\subsection{The common scalar flow}
Let $\mathcal H$ be a nonzero separable complex Hilbert space and let
$a$ be a positive elliptic, real-even homogeneous symbol of order
$m>1$. Put $s=m/2$ and
\[
 \rho_U=\|U\|_{\mathcal H}^2,\qquad
 M(U)=\int\rho_U,\qquad
 A(U)=\int a(\xi)\|\widehat U(\xi)\|_{\mathcal H}^2\,d\xi.
\]
As in the quantitative polarization theorem, take either
$\mathcal F(U)=\int\rho_U^\kappa$, $\kappa=1+m/N$, or the positive
definite critical Hartree interaction
$\mathcal F(U)=\mathcal D(\rho_U,\rho_U)$, $\kappa=2$.
Write $r=\kappa-1$. We allow a common real nonnegative multiplier
$b$ of order at most $m$, and write
\begin{equation}\label{eq:r12-internal-flow}
 iU_t=[a(D)+b(D)]U-V_{\rho_U}U,\qquad
 E_b(U)=\frac{A(U)+B_b(U)}2-\frac{\mathcal F(U)}{2\kappa}.
\end{equation}
No homogeneity or dilation inequality on $b$ is used in this section.

For the general flow statements we assume a conservative local
$H^s(\mathbb R^N;\mathcal H)$ theory with the continuation property:
every finite endpoint of the maximal interval entails an unbounded
$H^s$ norm. A concrete quartic class satisfying this assumption is
proved below. The sharp scalar mass $M_c$ is fixed by
\[
 \mathcal F(u)\le\kappa(M(u)/M_c)^r A(u).
\]

Define the positive trace-class operator on $\mathcal H$
\begin{equation}\label{eq:r12-Gamma}
 \Gamma(U)=\int |U(x)\rangle\langle U(x)|\,dx,
 \qquad \mu(U)=\|\Gamma(U)\|_{\mathrm{op}},
 \qquad \chi(U)=1-\frac{\mu(U)}{M(U)}.
\end{equation}
For a nonzero field, $\operatorname{Tr}\Gamma=M$ and $0\le\chi<1$.
The largest eigenvalue exists even if the internal space is
infinite-dimensional, because $\Gamma$ is compact.

\begin{lemma}[Conservation of all internal correlations]
\label{lem:r12-Gamma-conservation}
Every energy-space solution of \eqref{eq:r12-internal-flow} satisfies
$\Gamma(U(t))=\Gamma(U(0))$. In particular, $\mu$ and $\chi$ are
constant throughout its maximal interval.
\end{lemma}
\begin{proof}
For fixed internal vectors $h,k$, the scalar components
$u_h=\langle h,U\rangle$ and $u_k=\langle k,U\rangle$ solve the
same equation $iu_t=H(t)u$, where
$H(t)=a(D)+b(D)-V_{\rho_U}$ is a self-adjoint spatial form.
The Hilbert-triple product rule for
$C H^s\cap C^1 H^{-s}$ functions gives
\[
 \frac{d}{dt}\int u_h\overline{u_k}
 =-i\langle H(t)u_h,u_k\rangle
       +i\langle u_h,H(t)u_k\rangle=0.
\]
The local nonlinear form is continuous by H\"older at exponent
$2\kappa$; the Hartree form is continuous by HLS. Equality for a
countable dense set of $h,k$ proves equality of the operators.
The trace-norm estimate
\[
 \|\Gamma(U)-\Gamma(W)\|_1
 \le(\|U\|_2+\|W\|_2)\|U-W\|_2
\]
also shows that this identity passes to energy-space limits.
\end{proof}

The conserved eigenvalues measure the masses in mutually orthogonal
spatial modes after an internal unitary change of basis. Merely writing
one spatial profile in many components still gives a rank-one operator
and $\mu=M$.

\subsection{A strict improvement retaining the scalar sharp mass}
\begin{theorem}[Internal mass distribution and energy control]
\label{thm:r12-mixed-threshold}
There is a constant $c>0$, depending only on the scalar dispersion
and interaction and independent of the number of components, such that
every nonzero $U\in H^s(\mathcal H)$ satisfies
\begin{equation}\label{eq:r12-mixed-GN}
 \mathcal F(U)
 \le \kappa\left(\frac{M(U)}{M_c}\right)^r
        [1-c\chi(U)]A(U).
\end{equation}
One may take $c=\min\{C^{-1},1/2\}$, where $C$ is the
dimension-independent quantitative polarization constant.

Consequently, if the initial state satisfies
\begin{equation}\label{eq:r12-mixed-positive}
 \eta:=1-\left(\frac{M}{M_c}\right)^r(1-c\chi)>0,
\end{equation}
its conservative energy-space evolution is global in both time
directions, and
\begin{equation}\label{eq:r12-mixed-global-bound}
 \sup_{t\in\mathbb R}\{A(U(t))+B_b(U(t))\}
 \le\frac{2E_b(U(0))}{\eta}.
\end{equation}
Thus every fixed $\chi>0$ permits the mass interval
\[
 M<M_c(1-c\chi)^{-1/r},
\]
which extends strictly above the scalar critical mass.
\end{theorem}
\begin{proof}
Let $\mathfrak d$ and $\mathfrak p$ denote the normalized
interpolation deficit and polarization distance:
\[
 \mathfrak d(U)=1-
 \frac{\mathcal F(U)}{\kappa(M(U)/M_c)^r A(U)},
 \qquad
 \mathfrak p(U)=\inf_{\|\zeta\|=1,q\in H^s}
 \left\{\frac{M(U-\zeta q)}{M(U)}
              +\frac{A(U-\zeta q)}{A(U)}\right\}.
\]
For fixed $\zeta$, orthogonal projection gives
\[
 \inf_q M(U-\zeta q)=M-\langle\zeta,\Gamma(U)\zeta\rangle.
\]
Taking the infimum over unit $\zeta$ and discarding the nonnegative
kinetic term yields $\mathfrak p(U)\ge\chi(U)$.
The quantitative polarization theorem gives
$\mathfrak p(U)\le C\mathfrak d(U)$, proving
\eqref{eq:r12-mixed-GN}.

For clarity, if
$\Gamma_A=\int a(\xi)|\widehat U(\xi)\rangle
\langle\widehat U(\xi)|\,d\xi$, the exact distance formula is
\[
 \mathfrak p(U)=2-\lambda_{\max}
       \left(\frac{\Gamma(U)}M+\frac{\Gamma_A}A\right).
\]
The same scalar projection minimizes both terms in their combined
Hilbert norm. Since $\Gamma_A/A$ is positive with trace one, this
also gives $\mathfrak p\ge\chi$. No conservation of $\Gamma_A$ is
asserted.

By Lemma~\ref{lem:r12-Gamma-conservation}, $M,\chi$, and hence
$\eta$ are constant along the flow. Equation
\eqref{eq:r12-mixed-GN} gives
\[
 E_b(U(t))\ge\frac\eta2A(U(t))+\frac12B_b(U(t)).
\]
As $0<\eta\le1$, conservation yields
\eqref{eq:r12-mixed-global-bound}. Ellipticity, the conserved mass,
and the continuation property finish the proof.
\end{proof}

\begin{corollary}[Scalar reduction of critical-mass unbounded evolution]
\label{cor:r12-minimal-scalar}
Suppose $M(U)=M_c$ and a conservative orbit has unbounded kinetic
energy along a sequence of times, including possible approach to a
finite maximal endpoint. Then $\Gamma(U(0))$ has rank one, and there
is one fixed unit vector $\zeta\in\mathcal H$ such that
\[
 U(t,x)=\zeta u(t,x)
\]
throughout the maximal interval. Here $u$ solves the corresponding
scalar equation with mass $M_c$.

More generally, any energy-space-unbounded orbit of finite conserved
mass satisfies
\begin{equation}\label{eq:r12-nearcritical-mixing}
 \chi\le c^{-1}\left[1-(M_c/M)^r\right].
\end{equation}
\end{corollary}
\begin{proof}
If $M=M_c$ and $\chi>0$, Theorem~\ref{thm:r12-mixed-threshold}
applies with $\eta=c\chi$, contradicting unboundedness. Thus
$\mu=M$, so positivity and trace $M$ force
$\Gamma=M|\zeta\rangle\langle\zeta|$. For every
$h\perp\zeta$,
\[
 \int|\langle h,U(t,x)\rangle|^2\,dx
     =\langle h,\Gamma h\rangle=0.
\]
A countable orthonormal basis in $\zeta^\perp$ shows that
$U(t,x)$ lies in the same complex line almost everywhere.
Substitution gives the scalar equation. Finally an unbounded orbit
cannot satisfy \eqref{eq:r12-mixed-positive}; rearranging its
negation proves \eqref{eq:r12-nearcritical-mixing}.
\end{proof}

The scalar reduction holds for the initial datum and the entire
orbit, not merely for a rescaled concentration limit. For
$M=M_c(1+\varepsilon)$, \eqref{eq:r12-nearcritical-mixing} gives
$\chi\le(r/c)\varepsilon+O(\varepsilon^2)$.

\subsection{A concrete quartic flow}
\begin{proposition}[Unconditional realization in dimensions one to three]
\label{prop:r12-quartic-vector-flow}
Let $N=1,2,3$, let $a$ be any real-even positive elliptic homogeneous
quartic symbol, and let $b$ be real with
$0\le b(\xi)\le C(1+|\xi|^4)$. Then
\[
 iU_t=[a(D)+b(D)]U-\|U\|_{\mathcal H}^{8/N}U,
 \qquad U(0)\in H^2(\mathbb R^N;\mathcal H),
\]
has a conservative local $H^2$ flow with the continuation property.
Therefore Theorem~\ref{thm:r12-mixed-threshold} and
Corollary~\ref{cor:r12-minimal-scalar} hold without additional flow
assumptions for these equations, including countably many components.
\end{proposition}
\begin{proof}
The Fourier group is unitary on $H^2$. In these dimensions,
$H^2$ embeds into $L^\infty$, and the radial Hilbert-space map
$z\mapsto\|z\|^{8/N}z$ is locally Lipschitz on $H^2(\mathcal H)$.
Indeed its first three derivatives are bounded on bounded balls
because $8/N>2$; the chain rule, the $H^2$ product estimates, and
$H^1\hookrightarrow L^4$ give the required estimate, with constants
independent of internal dimension. Duhamel contraction gives the
local flow and its continuation property. Symmetric Fourier
Galerkin regularization gives mass and energy conservation; the
regularized solutions converge in $C H^2$ on every local existence
interval, and both conserved functionals are continuous there.
Lemma~\ref{lem:r12-Gamma-conservation} gives conservation of the
full internal operator.
\end{proof}

\paragraph{Relation to kinetic-energy inequalities.}
Conservation of all component overlaps and the use of the largest
occupation are established tools in nonlinear Schr\"odinger systems;
see Hong--Kwon--Yoon, arXiv:1712.06811, Section~1 and
Proposition~7.3. Their Theorem~2.3 also includes fractional
kinetic-energy inequalities. The specific connection used here is
the dimension-independent correction $1-c(1-\mu/M)$ to the
\emph{scalar sharp} critical-mass inequality. This gives a strict
extension above $M_c$ for every genuinely nonpolarized internal
state and the exact initial-data reduction at mass $M_c$.
The constant $c$ is not claimed to be optimal, and equality in
\eqref{eq:r12-mixed-positive} is not covered.
